% Discussion --- plan \S4.7.  Reordered from draft \S8 so that the chapter
% ends on its contribution rather than on its gaps.  The limitations
% subsection is new and gathers scope caveats previously scattered across
% the draft's Remark rem:flood-scope, \S7 and nowhere.

\section{Discussion}
\label{p:sec:discussion}

This chapter took the completed single-cell theory of \ChCore{} and
\ChDist{} and embedded it in a population. What came out is a renewal model
of within-host infection, forced by the intracellular process and exact in
its kernels; the demonstration that classical BMVR is a
projection of it, valid one exponential phase at a time; and an
effective-parameter map whose two ends differ by orders of magnitude. Three
consequences came with them: $R_0$ is invariant under the replacement of
budding by bursting, the flooding advantage at establishment is paid for in
invasion speed, and the construction is general enough to hold every release
mechanism of \cref{p:sec:reach} in one frame.

\subsection{Consequences for fitting practice}
\label{p:sec:fitting-consequences}

Three results of \cref{p:sec:projection} bear on how within-host models are
fitted, and each says what to do rather than only what is true.

A fitted $p$ is a property of the epidemic's phase, not of the pathogen, so
acute-phase and set-point estimates should differ systematically, in a
direction and by an amount \cref{p:eq:peff} specifies. Two published
estimates of $p$ for the same organism at different phases are therefore not
in conflict, and the ratio between them is itself an observable: by
\cref{p:eq:peff-span} it cannot exceed $\mean{X^2}_{\rm QS}$, so a reported
spread larger than that is evidence against the intracellular model rather
than against either fit.

The eclipse-like delay that within-host models insert by hand emerges here
from the intracellular dynamics with no extra parameter, its duration set by
the mean productive lifetime of \cref{p:eq:lifetime}. A model already
carrying a fitted eclipse and a bursting pathogen is thus fitting the same
delay twice, and the two estimates should be checked against each other
before either is trusted.

And no fit of $(p,d_{\Icell})$ recovers the three intracellular rates
(\cref{p:sec:identifiability}). The practical consequence is that a
population fit is worth doing only alongside one of the three single-cell
observables named there; with any one of them the system is
over-determined, and the fit becomes a test of the intracellular model
instead of a description of it. Which of the three is easiest to measure is
weighed in \ChDist{}.

\subsection{Limitations}
\label{p:sec:limitations}

The assumptions of \cref{p:sec:assumptions} are load-bearing, and each
would change something definite if relaxed. Holding target cells constant
(A1) makes the model linear, so supercritical solutions grow without bound
and the peak times of \cref{p:sec:overlays} exist only in the subcritical
cases; depletion would cap them, but the linear analysis remains the right
object for establishment and early growth, where the chapter's claims live.
The continuum convention (A2) is why $R_0=\gamma TV_\infty/c$ sits alongside
the branching mean $m=qV_\infty$: they agree when $\gamma T\ll c$ and are
otherwise different quantities, and confusing them would misread
\cref{p:thm:flood}. The synchronous initial condition (A3) is not innocuous,
since a non-trivial initial age distribution adds a term to
\cref{p:eq:Icell}. Single infection (A4) is a genuine restriction ---
multiply infected cells release superlinearly more, by
\cref{p:eq:gk,p:eq:Vk} --- and the multiplicity kernels exist but have not
been carried through the population layer. Absence of spatial structure (A5)
is the most consequential omission for a pathogen that spreads cell to cell.

Two further boundaries are worth naming. The branching comparison of
\cref{p:sec:flooding} holds the offspring mean fixed and is silent about
nonlinear effects; and the verification of \cref{p:app:verification} covers
reinfection only indirectly, through the characteristic equation.

The chapter is also written in dimensionless rates with $\lambda=1$, and
does not say where a real pathogen sits on \cref{p:fig:L-landscape}. For
\yp{} in the macrophage the interval is short and the biology resists a
clean measurement: the bacterium enters a permissive vacuole and replicates
there while the shift from the flea's ambient $26$--$28^\circ$C to the host's
$37^\circ$C induces the capsule, the pH~6 antigen and the type~III secretion
system, and then leaves by a route the primary literature still treats as
unresolved, with apoptosis, pyroptosis and necroptosis all reported. Placing
that organism on the $L=1$ boundary needs $\lambda$, $\mu$ and $\delta$ for
the intracellular phase --- and the exit mechanism is exactly what fixes
$\delta$.
% NEEDS-REF: Y. pestis macrophage residence, YCV maturation, and the
% YopJ/pyroptosis/necroptosis exit literature; see the two-type chapter's
% biology document for the account drawn on here

\subsection{Open problems}
\label{p:sec:open-problems}

Four questions remain open on the population side.

\begin{enumerate}
\item \textbf{Two-type killed second moments, and the two-type $D$.} The
release kernel of a two-type renewal model needs
$\mean{X_\alpha^2\ind\{\tau_c>\alpha\}}$,
$\mean{X_\alpha Y_\alpha\ind\{\tau_c>\alpha\}}$,
$\mean{Y_\alpha^2\ind\{\tau_c>\alpha\}}$, and the analogue of $D(t)$,
extinction without catastrophe. The moment hierarchy does not close upward,
exactly as in one type; the one-type escape was to differentiate the Riccati
equation, and the two-type route should be to differentiate the backward
system of \ChTwoType{}. This is the main technical obstacle between this
chapter and an eclipse-aware renewal model, and it should not be assumed
easy.
\item \textbf{The flooding boundary against real parameters.} Map
$\delta(\lambda+\mu)>\lambda\mu$ for \yp{} in the macrophage, and compute
the \emph{size} of the advantage from \cref{p:eq:flood}: does the pathogen
sit on the advantageous side, and by how much? The estimates discussed in \ChDist{} --- the near-certain bursting of \ft{}
\cite{carruthers2020stochastic}, the clearance-dominated anthrax case ---
suggest that real pathogens sit on both sides of the boundary, which is what
makes the question worth asking.
\item \textbf{Population-level variance.} Propagate $\mathrm{Var}(W)$ of
\cref{p:eq:EW2} through the branching structure to a whole-host prediction.
For single-cell assays the first moments suffice; for whole-host data the
variance is the main stochastic signal that classical BMVR cannot
reproduce. The integer-valued nature of the offspring count, whose
distribution matters for establishment
\cite{williams2024reproduction}, enters through the same quantity.
\item \textbf{Partial release and the flooding boundary in $\varphi$.}
Where on the continuum of \cref{p:sec:partial} does the flooding advantage
survive, and does the closed-form structure persist at intermediate
$\varphi$?
\end{enumerate}

Two smaller questions are recorded in the text where they arise, because
each is now a single definite statement rather than an open direction. The
hypothesis of \cref{p:prop:peff-monotone} --- that $\rho=g/S$ increases with
cell age --- is a claim about the one closed-form function
\cref{p:eq:rho-closed} of one variable, so proving it is finite algebra; and
the ordering $r_{\rm bud}>r_{\rm burst}$ is by \cref{p:prop:rorder}
equivalent to the single Laplace inequality \cref{p:eq:laplace-order}, with
the coupling constants eliminated. Both are supported by wide numerical
sweeps (tests~M and~N of \cref{p:app:verification}) and neither is proved.
The single-cell open problems belong to \ChDist{}.

\subsection{Connections forward}
\label{p:sec:forward}

The two-type process of \ChTwoType{} supplies the most immediate
continuation, and it is where the chapter's principal technical obstacle
lives. In the GATE regime $\delta_1=0<\delta_2$ --- no burst before
conversion --- the no-catastrophe probability has zero initial hazard and
then rises: precisely the eclipse-phase signature that within-host models
insert by hand, emerging mechanistically, with the eclipse duration set by
the conversion rate. The construction here is ready to receive the two-type
kernels the moment they exist; nothing in \cref{p:eq:skeleton} need change
to accept them. The bursting--budding comparison finds its evolutionary home
in \ChPath{}, where the intermediate mechanisms of \cref{p:sec:reach} supply
the phenotypes a selection argument needs. The nonlinear effects that
\fwd{nonlinear}{the later modelling chapters} treat --- immune saturation,
local depletion, the all-run-out-at-once intuition --- are exactly what the
branching comparison deliberately excludes, and should be named wherever the
flooding advantage is invoked.

\FloatBarrier
